% SPDX-License-Identifier: GPL-3.0-or-later
\chapter{Baremul Complet al celor 74 de Întrebări Quiz și Glosar Astronomic}
\label{chap:barem-glosar}

Acest capitol centralizează răspunsurile corecte și explicațiile științifice pentru toate cele \textbf{74 de întrebări grilă de verificare} din capitolele anterioare, urmate de un \textbf{Glosar de Termeni Astronomici și Spațiali}.

\section{Baremul celor 74 de Întrebări Grilă de Verificare}

\begingroup
\small
\begin{longtable}{@{}r >{\raggedright\arraybackslash}p{2.7cm} >{\raggedright\arraybackslash}p{4.8cm} >{\raggedright\arraybackslash}p{6.5cm}@{}}
\caption{Baremul complet de verificare pentru cele 74 de corpuri cerești, misiuni, exoplanete și galaxii} \\
\toprule
\textbf{\#} & \textbf{Obiect Cosmic} & \textbf{Întrebarea Quiz} & \textbf{Răspuns Corect \& Explicație} \\
\midrule
\endfirsthead
\multicolumn{4}{c}{{\footnotesize\itshape (continuare barem Quiz de pe pagina anterioară)}} \\
\toprule
\textbf{\#} & \textbf{Obiect Cosmic} & \textbf{Întrebarea Quiz} & \textbf{Răspuns Corect \& Explicație} \\
\midrule
\endhead
\bottomrule
\endfoot
\textbf{1} & \textbf{Soarele} & Ce procent din masa întregului Sistem Solar se află în Soare? & \textbf{C: Peste 99,8\%}. {\footnotesize Soarele concentrează 99,86\% din masa Sistemului Solar; toate planetele, lunile și asteroizii la un loc reprezintă doar 0,14\%!} \\ \midrule
\textbf{2} & \textbf{Mercur} & Cât durează un an pe planeta Mercur (o orbită completă în jurul Soarelui)? & \textbf{B: 88 de zile pământești}. {\footnotesize Fiind cea mai apropiată planetă de Soare, Mercur parcurge orbita cu aproximativ 47 kilometri pe secundă și încheie un an în doar 88 de zile terestre.} \\ \midrule
\textbf{3} & \textbf{Venus} & De ce este Venus mai fierbinte decât Mercur, deși se află mai departe de Soare? & \textbf{B: Are o atmosferă densă de CO$_2$ care creează un efect de seră extrem}. {\footnotesize Atmosfera groasă de dioxid de carbon captează radiația termică, ridicând temperatura la +464 $^\circ$C pe toată suprafața planetei.} \\ \midrule
\textbf{4} & \textbf{Pământul} & Ce ne protejează planeta de particulele încărcate electric ale vântului solar? & \textbf{A: Câmpul magnetic (magnetosfera) generat de nucleul Pământului}. {\footnotesize Nucleul extern din fier lichid al Pământului funcționează ca un dinam și creează magnetosfera, care deviază vântul solar.} \\ \midrule
\textbf{5} & \textbf{Luna} & De ce vedem de pe Pământ întotdeauna aceeași față a Lunii? & \textbf{B: Pentru că perioada de rotație a Lunii este egală cu perioada ei de orbitare în jurul Pământului}. {\footnotesize Fenomenul se numește rotație sincronă: Luna se rotește în jurul axei sale în exact 27,3 zile, adică în același timp în care înconjoară Pământul.} \\ \midrule
\textbf{6} & \textbf{Marte} & Cum se numește cel mai înalt vulcan de pe Marte și din întregul Sistem Solar? & \textbf{B: Olympus Mons}. {\footnotesize Olympus Mons are o altitudine de aproape 22 de kilometri și o bază de aproximativ 600 de kilometri diametru!} \\ \midrule
\textbf{7} & \textbf{Phobos} & De ce răsare Phobos din vest și apune în est pe cerul planetei Marte? & \textbf{A: Pentru că se rotește în jurul lui Marte în 7 ore și 39 de minute, mai repede decât se învârte Marte în jurul axei sale}. {\footnotesize Phobos înconjoară Marte în doar 7,65 ore, în timp ce o zi pe Marte durează 24,6 ore, deci satelitul depășește rotația planetei.} \\ \midrule
\textbf{8} & \textbf{Deimos} & Cum arată suprafața satelitului Deimos în comparație cu cea a lui Phobos? & \textbf{B: Este mult mai netedă, deoarece craterele sunt umplute cu un strat gros de praf (regolit)}. {\footnotesize Deimos este acoperit cu un strat gros de regolit pulverulent care a umplut craterele vechi, dându-i un aspect catifelat.} \\ \midrule
\textbf{9} & \textbf{Jupiter} & Ce este celebra \guillemotleft{}Mare Pată Roșie\guillemotright{} de pe Jupiter? & \textbf{C: O furtună anticiclonică mai mare decât Pământul}. {\footnotesize Marea Pată Roșie este un anticiclon gigantic cu vânturi de peste 430 kilometri pe oră, observat de mai bine de 300 de ani.} \\ \midrule
\textbf{10} & \textbf{Io} & Ce fenomen provoacă activitatea vulcanică extremă de pe satelitul Io? & \textbf{B: Încălzirea mareică produsă de gravitația lui Jupiter și a lunilor vecine}. {\footnotesize Forțele mareice uriașe ale lui Jupiter deformează continuu interiorul lui Io, generând căldură prin frecare internă.} \\ \midrule
\textbf{11} & \textbf{Europa} & De ce sunt oamenii de știință atât de interesați de satelitul Europa? & \textbf{B: Adăpostește sub gheață un ocean global de apă lichidă care ar putea susține viață}. {\footnotesize Oceanul subglaciar al Europei, aflat în contact cu un fund oceanic stâncos, oferă ingredientele esențiale pentru potențiale ecosisteme.} \\ \midrule
\textbf{12} & \textbf{Ganymede} & Prin ce este Ganymede unic printre toți sateliții naturali din Sistemul Solar? & \textbf{A: Este cel mai mare satelit (mai mare decât Mercur) și are câmp magnetic propriu}. {\footnotesize Ganymede depășește în diametru planeta Mercur și este singurul satelit cu o magnetosferă proprie.} \\ \midrule
\textbf{13} & \textbf{Callisto} & Cum se numește uriașul bazin de impact cu inele concentrice de pe Callisto? & \textbf{A: Valhalla}. {\footnotesize Bazinul Valhalla este cel mai mare crater multi-inelar de pe Callisto, având un diametru de aproape 3.800 de kilometri.} \\ \midrule
\textbf{14} & \textbf{Saturn} & Din ce sunt alcătuite în principal celebrele inele ale lui Saturn? & \textbf{C: Din miliarde de fragmente de gheață și rocă}. {\footnotesize Inelele lui Saturn nu sunt solide, ci sunt formate din nenumărate particule și blocuri de gheață care orbitează independent planeta.} \\ \midrule
\textbf{15} & \textbf{Mimas} & Ce formațiune geologică celebră ocupă o treime din diametrul satelitului Mimas? & \textbf{A: Craterul Herschel}. {\footnotesize Craterul Herschel are 139 de kilometri lățime și un munte central înalt de 6 kilometri, dându-i lui Mimas aspectul său inconfundabil.} \\ \midrule
\textbf{16} & \textbf{Enceladus} & Ce alimentează inelul E al planetei Saturn? & \textbf{B: Gheizerele de gheață și vapori de apă de la polul sud al lui Enceladus}. {\footnotesize Jetul criovulcanic al lui Enceladus expulzează continuu granule de gheață care scapă de gravitația micului satelit și formează inelul E.} \\ \midrule
\textbf{17} & \textbf{Tethys} & Din ce material este alcătuită în proporție de peste 93\% luna Tethys a lui Saturn? & \textbf{B: Din gheață de apă aproape pură (având o densitate mai mică decât a apei lichide)}. {\footnotesize Cu o densitate medie de 0,984 g/cm$^3$, Tethys este alcătuită aproape exclusiv din gheață de apă, reflectând puternic lumina Soarelui.} \\ \midrule
\textbf{18} & \textbf{Dione} & Ce s-au dovedit a fi în realitate dungile albe luminoase (\guillemotleft{}wispy terrain\guillemotright{}) de pe satelitul Dione, fotografiate de aproape de sonda Cassini? & \textbf{B: Faleze și canioane tectonice abrupte din gheață strălucitoare (precum Padua Chasmata)}. {\footnotesize Imaginile de înaltă rezoluție de la Cassini au arătat că dungile luminoase sunt pereți abrupți de falii tectonice unde gheața curată este expusă la suprafață.} \\ \midrule
\textbf{19} & \textbf{Rhea} & Ce loc ocupă Rhea ca mărime între sateliții planetei Saturn? & \textbf{B: Este al doilea satelit ca mărime al lui Saturn, după Titan}. {\footnotesize Cu o rază de 763,8 kilometri (diametru de 1.528 de kilometri), Rhea este a doua lună ca mărime a lui Saturn, fiind depășită doar de Titan.} \\ \midrule
\textbf{20} & \textbf{Titan} & Din ce sunt formate lacurile și mările de la suprafața lui Titan? & \textbf{B: Din metan și etan lichid}. {\footnotesize La temperatura de -179 $^\circ$C de pe Titan, apa este tare ca piatra, iar metanul și etanul formează nori, ploi, râuri și mări.} \\ \midrule
\textbf{21} & \textbf{Hyperion} & De ce seamănă satelitul Hyperion al lui Saturn cu un burete sau un fagure înghețat? & \textbf{A: Are o porozitate internă de peste 40\%, iar impacturile comprimă gheața în cratere adânci cu margini ascuțite}. {\footnotesize Cu o densitate de doar 0,54 g/cm$^3$, Hyperion este extrem de poros, astfel încât impacturile meteoritice creează cavități adânci prin compresie, fără a le șterge pe cele vecine.} \\ \midrule
\textbf{22} & \textbf{Iapetus} & Prin ce se remarcă vizual satelitul Iapetus al lui Saturn? & \textbf{A: Are o emisferă neagră ca cărbunele și una albă ca zăpada, plus o creastă muntoasă la ecuator}. {\footnotesize Iapetus prezintă cel mai dramatic contrast de culoare (albedo) din Sistemul Solar și o creastă ecuatorială care îi dă forma unei nuci.} \\ \midrule
\textbf{23} & \textbf{Phoebe} & Care este originea satelitului exterior Phoebe al lui Saturn și ce efect produce praful desprins de pe el? & \textbf{A: Este un corp capturat din Centura Kuiper, iar praful său întunecat colorează în negru o emisferă a satelitului Iapetus}. {\footnotesize Orbita retrogradă și compoziția lui Phoebe confirmă că a fost capturat din Centura Kuiper, iar praful din Inelul Phoebe este responsabil pentru emisfera întunecată a lui Iapetus.} \\ \midrule
\textbf{24} & \textbf{Uranus} & Ce gaz din atmosfera lui Uranus îi oferă culoarea turcoaz (cyan)? & \textbf{C: Metanul}. {\footnotesize Metanul (CH$_4$) din atmosfera superioară absoarbe componenta roșie a luminii solare și difuzează lumina albastră-verzuie.} \\ \midrule
\textbf{25} & \textbf{Miranda} & Ce record deține stânca Verona Rupes de pe satelitul Miranda al lui Uranus? & \textbf{B: Este cea mai înaltă faleză (prăpastie) cunoscută din Sistemul Solar (aproximativ 20 de kilometri înălțime)}. {\footnotesize Verona Rupes se ridică la aproximativ 20 de kilometri deasupra terenului înconjurător, fiind de peste două ori mai înaltă decât muntele Everest.} \\ \midrule
\textbf{26} & \textbf{Ariel} & Prin ce se distinge suprafața satelitului Ariel față de vecina sa întunecată Umbriel? & \textbf{A: Ariel este cea mai luminoasă lună a lui Uranus și are o suprafață tânără, tăiată de canioane tectonice (Kachina Chasmata) și câmpii criovulcanice}. {\footnotesize Încălzirea mareică din trecut a remodelat suprafața lui Ariel prin falii tectonice și criovulcanism, făcând-o cea mai strălucitoare lună a lui Uranus.} \\ \midrule
\textbf{27} & \textbf{Umbriel} & Ce detaliu strălucitor iese în evidență pe suprafața foarte întunecată a satelitului Umbriel? & \textbf{A: Inelul alb luminos de pe podeaua craterului ecuatorial Wunda}. {\footnotesize Craterul Wunda conține un inel lat și foarte luminos (probabil gheață de dioxid de carbon sau gheață de apă curată expusă la impact), care contrastează puternic cu restul suprafeței întunecate a lui Umbriel.} \\ \midrule
\textbf{28} & \textbf{Titania} & De unde provin numele sateliților planetei Uranus, precum Titania și Miranda? & \textbf{B: Din operele literare ale lui William Shakespeare și Alexander Pope}. {\footnotesize Spre deosebire de celelalte planete ale căror luni poartă nume din mitologia greco-romană, lunile lui Uranus sunt numite după personaje din Shakespeare și Pope.} \\ \midrule
\textbf{29} & \textbf{Oberon} & Ce formațiune spectaculoasă a observat sonda Voyager 2 ridicându-se pe marginea orizontului (limbul) satelitului Oberon? & \textbf{A: Un munte de aproximativ 11 kilometri înălțime (probabil piscul central al unui mare bazin de impact)}. {\footnotesize În imaginile Voyager 2, pe limbul sud-estic al lui Oberon se profilează un munte înalt de aproximativ 11 kilometri, mai înalt decât Everestul de pe Pământ.} \\ \midrule
\textbf{30} & \textbf{Neptun} & Cum a fost descoperită planeta Neptun în anul 1846? & \textbf{B: Prin calcule matematice bazate pe perturbațiile orbitei lui Uranus}. {\footnotesize Matematicienii au observat că Uranus nu urma exact orbita calculată și au dedus poziția unei planete necunoscute care îl atrăgea gravitațional.} \\ \midrule
\textbf{31} & \textbf{Triton} & Ce ne indică faptul că satelitul Triton se rotește în sens invers (retrograd) în jurul lui Neptun? & \textbf{A: Că nu s-a format lângă Neptun, ci a fost un corp din Centura Kuiper capturat de gravitația planetei}. {\footnotesize Un satelit format împreună cu planeta sa orbitează în sensul rotației planetei; orbita retrogradă a lui Triton dovedește că a fost capturat din Centura Kuiper.} \\ \midrule
\textbf{32} & \textbf{Ceres} & Unde se află planeta pitică Ceres? & \textbf{B: În Centura Principală de Asteroizi, între Marte și Jupiter}. {\footnotesize Ceres orbitează la 2,77 UA de Soare, în mijlocul Centurii Principale de Asteroizi dintre Marte și Jupiter.} \\ \midrule
\textbf{33} & \textbf{Pluto} & Cum se numește celebra regiune strălucitoare în formă de inimă de pe Pluto? & \textbf{A: Tombaugh Regio (Sputnik Planitia)}. {\footnotesize Regiunea în formă de inimă a fost numită Tombaugh Regio în onoarea descoperitorului lui Pluto, iar lobul ei vestic este ghețarul Sputnik Planitia.} \\ \midrule
\textbf{34} & \textbf{Charon} & De ce stau Pluto și Charon mereu cu aceeași față îndreptată unul spre celălalt? & \textbf{A: Ambele corpuri sunt în rotație sincronă reciprocă (blocare mareică dublă)}. {\footnotesize Pluto și Charon se rotesc în jurul axelor lor în exact 6,387 zile, având exact aceeași durată în care orbitează în jurul baricentrului lor comun.} \\ \midrule
\textbf{35} & \textbf{Haumea} & De ce are planeta pitică Haumea o formă alungită, asemănătoare unei mingi de rugby? & \textbf{B: Din cauza rotației extrem de rapide (o rotație completă în mai puțin de 4 ore)}. {\footnotesize Forța centrifugă produsă de rotația în doar 3,9 ore i-a deformat corpul într-un elipsoid alungit.} \\ \midrule
\textbf{36} & \textbf{Makemake} & Care a fost porecla dată de astronomi planetei pitice Makemake la descoperire? & \textbf{B: Iepurașul de Paște (Easterbunny)}. {\footnotesize Fiind descoperită la scurt timp după Paștele anului 2005, echipa a poreclit-o \guillemotleft{}Easterbunny\guillemotright{}, iar numele oficial Makemake provine din mitologia Rapa Nui (Insula Paștelui).} \\ \midrule
\textbf{37} & \textbf{Eris} & Ce eveniment astronomic a determinat descoperirea planetei pitice Eris în 2005? & \textbf{B: Definirea oficială a categoriei de \guillemotleft{}planetă pitică\guillemotright{} și reclasificarea lui Pluto în 2006}. {\footnotesize Deoarece Eris s-a dovedit mai masivă decât Pluto, Uniunea Astronomică Internațională a stabilit în 2006 criteriile pentru planete și planete pitice.} \\ \midrule
\textbf{38} & \textbf{4 Vesta} & Ce au descoperit oamenii de știință în centrul bazinului Rheasilvia de pe asteroidul Vesta? & \textbf{B: Un munte central înalt de 22 de kilometri, aproape cât Olympus Mons de pe Marte}. {\footnotesize Impactul colosal care a creat craterul Rheasilvia a ridicat în centru un munte de 22 de kilometri înălțime, unul dintre cele mai înalte vârfuri din Sistemul Solar.} \\ \midrule
\textbf{39} & \textbf{2 Pallas} & De ce seamănă suprafața asteroidului Pallas cu o minge de golf? & \textbf{B: Orbita sa foarte înclinată (34,8$^\circ$) produce ciocniri extrem de violente și numeroase cu alți asteroizi}. {\footnotesize Traversând centura de asteroizi sub un unghi mare, Pallas suferă impacturi la viteze de până la patru ori mai mari decât media din centură.} \\ \midrule
\textbf{40} & \textbf{16 Psyche} & Din ce este alcătuit în proporție ridicată asteroidul 16 Psyche? & \textbf{B: Din fier și nichel metalic (și silicați)}. {\footnotesize 16 Psyche este cel mai mare asteroid de tip M, bogat în fier și nichel, oferind indicii despre cum s-au format nucleele planetelor telurice.} \\ \midrule
\textbf{41} & \textbf{101955 Bennu} & Ce misiune spațială a colectat mostre de pe asteroidul Bennu și le-a adus pe Pământ în 2023? & \textbf{B: OSIRIS-REx}. {\footnotesize Sonda NASA OSIRIS-REx a colectat peste 120 de grame de regolit carbonic de pe Bennu și a parașutat capsula în deșertul Utah în septembrie 2023.} \\ \midrule
\textbf{42} & \textbf{99942 Apophis} & Ce se va întâmpla pe 13 aprilie 2029 cu asteroidul Apophis? & \textbf{B: Va trece în deplină siguranță la aproximativ 32.000 de kilometri de Pământ, fiind vizibil cu ochiul liber}. {\footnotesize Calculele radar au demonstrat că Apophis va trece fără niciun pericol la aproximativ 32.000 de kilometri de Pământ, oferind astronomilor ocazia rară de a-l studia de aproape.} \\ \midrule
\textbf{43} & \textbf{Stația Spațială Internațională (ISS)} & În cât timp face Stația Spațială Internațională o tură completă în jurul Pământului? & \textbf{B: În aproximativ 92 de minute}. {\footnotesize La viteza orbitală de aproximativ 28.000 kilometri pe oră, ISS înconjoară întregul Pământ în aproximativ 92 de minute ( de aproape 16 ori pe zi).} \\ \midrule
\textbf{44} & \textbf{Telescopul Spațial Hubble} & De ce vede Telescopul Hubble mult mai clar decât multe telescoape de pe Pământ? & \textbf{B: Pentru că se află în spațiu, deasupra atmosferei turbulente a Pământului}. {\footnotesize Atmosfera Pământului estompează și face stelele să \guillemotleft{}pâlpâie\guillemotright{} (scintilație); aflat în vidul spațial la 535 de kilometri altitudine, Hubble obține imagini de o claritate excepțională.} \\ \midrule
\textbf{45} & \textbf{Telescopul Spațial James Webb (JWST)} & Din câte segmente hexagonale aurite este alcătuită oglinda principală a Telescopului James Webb? & \textbf{B: 18 segmente}. {\footnotesize Oglinda de 6,5 metri a JWST este formată din 18 segmente hexagonale din beriliu aurit care s-au depliat și aliniat la nivel nanometric în spațiu.} \\ \midrule
\textbf{46} & \textbf{Telescopul Spațial Euclid} & Care este obiectivul științific principal al misiunii europene Euclid? & \textbf{B: Realizarea unei hărți 3D a miliarde de galaxii pentru a studia materia și energia întunecată}. {\footnotesize Euclid măsoară geometria la scară mare a Universului pentru a înțelege cum materia întunecată și energia întunecată îi modelează expansiunea.} \\ \midrule
\textbf{47} & \textbf{Voyager 1} & Ce poartă la bord sondele Voyager 1 și Voyager 2 ca mesaj pentru eventuale civilizații extraterestre? & \textbf{A: Discul de Aur (Golden Record) cu sunete, imagini și saluturi în 55 de limbi, inclusiv în română}. {\footnotesize Discul de Aur Voyager conține 115 imagini, muzică de pe Pământ și saluturi în 55 de limbi (inclusiv \guillemotleft{}Salutări la toată lumea\guillemotright{} în limba română).} \\ \midrule
\textbf{48} & \textbf{Voyager 2} & Care este singura sondă spațială care a vizitat planetele Uranus și Neptun? & \textbf{B: Voyager 2}. {\footnotesize Voyager 2 este singura navă spațială care a survolat toate cele patru planete gigante, inclusiv Uranus (1986) și Neptun (1989).} \\ \midrule
\textbf{49} & \textbf{New Horizons} & În cât timp a ajuns sonda New Horizons de la Pământ până la orbita Lunii după lansarea din 2006? & \textbf{B: În doar 9 ore}. {\footnotesize Lansată cu peste 58.000 kilometri pe oră, New Horizons a traversat distanța Pământ--Lună în doar 9 ore!} \\ \midrule
\textbf{50} & \textbf{Parker Solar Probe} & Ce record absolut deține sonda Parker Solar Probe? & \textbf{B: Cel mai rapid obiect construit vreodată de om (aproximativ 690.000 kilometri pe oră) și cel mai apropiat de Soare}. {\footnotesize Atrăgând-o gravitația imensă a Soarelui la periheliu, Parker Solar Probe atinge aproximativ 191 kilometri pe secundă (690.000 kilometri pe oră).} \\ \midrule
\textbf{51} & \textbf{Rosetta \& Philae} & Ce premieră mondială a realizat misiunea europeană Rosetta în anul 2014? & \textbf{B: A intrat pe orbita unei comete (67P) și a așezat robotul Philae pe suprafața ei}. {\footnotesize Rosetta și modulul Philae au realizat prima orbitare și prima asolizare din istorie pe nucleul unei comete.} \\ \midrule
\textbf{52} & \textbf{JUICE (Jupiter Icy Moons Explorer)} & În jurul cărui satelit natural va intra pe orbită sonda europeană JUICE în anul 2034? & \textbf{B: Ganymede (cel mai mare satelit al lui Jupiter)}. {\footnotesize JUICE va deveni prima sondă din istorie care va intra pe orbită în jurul unui satelit al altei planete: Ganymede.} \\ \midrule
\textbf{53} & \textbf{Perseverance \& Ingenuity} & Ce a demonstrat în premieră instrumentul MOXIE de la bordul roverului Perseverance? & \textbf{B: Că se poate produce oxigen pur din atmosfera de dioxid de carbon a lui Marte}. {\footnotesize MOXIE a extras oxigen din CO$_2$-ul atmosferic marțian, o tehnologie esențială pentru viitorii astronauți (pentru respirație și combustibil de rachetă).} \\ \midrule
\textbf{54} & \textbf{Curiosity (Mars Science Laboratory)} & Ce descoperire fundamentală a făcut roverul Curiosity în craterul Gale de pe Marte? & \textbf{B: Că în trecutul îndepărtat al lui Marte au existat lacuri și râuri de apă lichidă cu condiții prielnice vieții microbiene}. {\footnotesize Analizând pietrișul rotunjit de apă și argilele din craterul Gale, Curiosity a demonstrat că Marte a găzduit lacuri străvechi de apă potabilă.} \\ \midrule
\textbf{55} & \textbf{Norul lui Oort} & Ce se găsește în principal în Norul lui Oort, la aproximativ 1 an-lumină de Soare? & \textbf{A: Miliarde de nuclee înghețate de comete la granița Sistemului Solar}. {\footnotesize Norul lui Oort este rezervorul sferic uriaș de comete înghețate aflat la limita atracției gravitaționale a Soarelui (aproximativ 1 an-lumină).} \\ \midrule
\textbf{56} & \textbf{Proxima Centauri \& Proxima b} & La ce distanță de noi se află Proxima Centauri, cea mai apropiată stea de Sistemul Solar? & \textbf{A: 4,24 ani-lumină}. {\footnotesize Proxima Centauri se află la 4,24 ani-lumină (aproximativ 40.000 de miliarde de kilometri) de Sistemul nostru Solar.} \\ \midrule
\textbf{57} & \textbf{Sirius A \& B (Luceafărul Câinelui)} & Ce tip de stea este companionul Sirius B care orbitează în jurul strălucitoarei Sirius A? & \textbf{A: O pitică albă extrem de densă, de mărimea Pământului dar cu masa Soarelui}. {\footnotesize Sirius B este o pitică albă, reprezentând nucleul dens rămas dintr-o stea străveche, având masa Soarelui strânsă într-un volum cât Pământul.} \\ \midrule
\textbf{58} & \textbf{Sistemul TRAPPIST-1 (7 Exoplanete)} & Câte planete stâncoase de mărimea Pământului orbitează în jurul stelei TRAPPIST-1? & \textbf{A: 7 exoplanete stâncoase (dintre care 3 în zona locuibilă)}. {\footnotesize Sistemul TRAPPIST-1 deține recordul pentru cele mai multe planete telurice de mărimea Pământului (7 planete: b, c, d, e, f, g, h) în jurul unei singure stele.} \\ \midrule
\textbf{59} & \textbf{Exoplaneta 55 Cancri e (Janssen)} & De ce este suprafața exoplanetei 55 Cancri e acoperită de un ocean de lavă topită? & \textbf{A: Pentru că orbitează extrem de aproape de steaua ei (un an durează doar aproximativ 18 ore), la aproximativ 2.000 $^\circ$C}. {\footnotesize Fiind la doar 0,015 UA de steaua ei, 55 Cancri e este încălzită la aproximativ 2.000 $^\circ$C, topind rocile de la suprafață într-un ocean global de magmă.} \\ \midrule
\textbf{60} & \textbf{Roiul Stelar Pleiadele (\quotedblbase{}Cloșca cu Pui\textquotedblright{})} & Cum este numit roiul stelar Pleiadele (M45) în tradiția populară românească și acum câți ani a plecat lumina lui spre noi? & \textbf{A: \quotedblbase{}Cloșca cu Pui\textquotedblright{}, iar lumina a plecat acum aproximativ 444 de ani}. {\footnotesize În tradiția românească, Pleiadele sunt numite \quotedblbase{}Cloșca cu Pui\textquotedblright{} (sau \quotedblbase{}Găinușa\textquotedblright{}), iar la distanța de 444 de ani-lumină, lumina lor a plecat spre Pământ în secolul al XVI-lea.} \\ \midrule
\textbf{61} & \textbf{Betelgeuse (Supergiganta Roșie)} & Până unde s-ar întinde supergiganta roșie Betelgeuse dacă ar fi așezată în locul Soarelui nostru? & \textbf{A: Ar înghiți Mercur, Venus, Pământul și Marte, ajungând aproape de orbita lui Jupiter}. {\footnotesize Având o rază de aproximativ 700 de ori mai mare decât a Soarelui (aproximativ 3,5 UA), Betelgeuse ar depăși orbita lui Marte și Centura de Asteroizi, apropiindu-se de Jupiter.} \\ \midrule
\textbf{62} & \textbf{Nebuloasa Orion (M42)} & Ce se întâmplă în interiorul Nebuloasei Orion (M42), studiată de Telescopul James Webb? & \textbf{A: Se nasc sute de stele noi și sisteme planetare tinere din nori de gaz și praf cosmic}. {\footnotesize Nebuloasa Orion este cea mai apropiată regiune masivă de formare a stelelor (\quotedblbase{}maternitate stelară\textquotedblright{}), unde JWST observă stele nou-născute și discuri protoplanetare.} \\ \midrule
\textbf{63} & \textbf{Stâlpii Creației (Nebuloasa Vulturului M16)} & Ce a reușit să vadă Telescopul Spațial James Webb în interiorul Stâlpilor Creației folosind lumina infraroșie? & \textbf{A: A privit prin perdeaua de praf și a dezvăluit stelele tinere roșii care se nasc în vârful coloanelor}. {\footnotesize Lumina infraroșie pătrunde prin norii opaci de praf cosmic, permițând telescopului James Webb să fotografieze protostelele roșii nou-născute din interiorul Stâlpilor Creației.} \\ \midrule
\textbf{64} & \textbf{Sagittarius A* (Gaura Neagră Supermasivă)} & Ce este \quotedblbase{}orizontul evenimentelor\textquotedblright{} al găurii negre Sagittarius A* din centrul Căii Lactee? & \textbf{A: Limita în jurul găurii negre dincolo de care gravitația este atât de puternică încât nici lumina nu mai poate ieși}. {\footnotesize Orizontul evenimentelor este granița gravitațională a unei găuri negre unde viteza necesară pentru a scăpa depășește viteza luminii (300.000 kilometri pe secundă).} \\ \midrule
\textbf{65} & \textbf{Calea Lactee (Galaxia Noastră)} & Unde se află Sistemul nostru Solar în cadrul Galaxiei Calea Lactee? & \textbf{A: În Brațul Orion, la aproximativ 26.000 de ani-lumină de centrul galaxiei}. {\footnotesize Soarele și planetele noastre se află în Brațul Orion, la aproximativ 26.000 de ani-lumină de centrul Căii Lactee (cam la jumătatea razei discului galactic).} \\ \midrule
\textbf{66} & \textbf{Galaxia Andromeda (M31)} & Care este cel mai îndepărtat obiect cosmic vizibil cu ochiul liber pe cerul nopții și la ce distanță se află? & \textbf{A: Galaxia Andromeda (M31), aflată la aproximativ 2,5 milioane de ani-lumină}. {\footnotesize Galaxia Andromeda (M31), situată la 2,54 de milioane de ani-lumină, poate fi văzută cu ochiul liber într-o noapte fără poluare luminoasă.} \\ \midrule
\textbf{67} & \textbf{Nebuloasa Carina (\quotedblbase{}Stâncile Cosmice\textquotedblright{}, NGC 3372)} & Cum se numește formațiunea spectaculoasă din Nebuloasa Carina fotografiată în prima serie de imagini a Telescopului James Webb? & \textbf{A: Stâncile Cosmice, creste de gaz și praf înalte de 7 ani-lumină unde se nasc stele}. {\footnotesize \quotedblbase{}Stâncile Cosmice\textquotedblright{} din NGC 3324 (Nebuloasa Carina) reprezintă marginea unei cavități gazoase uriașe sculptate de stele tinere masive.} \\ \midrule
\textbf{68} & \textbf{Nebuloasa Crabului (M1)} & Ce obiect extraordinar se află chiar în centrul Nebuloasei Crabului (M1)? & \textbf{A: Un pulsar (o stea neutronică densă cât un oraș) care se rotește de 30 de ori pe secundă}. {\footnotesize În centrul Nebuloasei Crabului se află Pulsarul din Crab, o stea neutronică rămasă după supernova din anul 1054.} \\ \midrule
\textbf{69} & \textbf{Nebuloasa Inelul (M57)} & Cum ia naștere o \quotedblbase{}nebuloasă planetară\textquotedblright{} precum Nebuloasa Inelul (M57)? & \textbf{A: O stea de mărime medie (ca Soarele) își eliberează straturile exterioare de gaz la finalul vieții, lăsând în centru o pitică albă}. {\footnotesize Nebuloasele planetare sunt învelișuri strălucitoare de gaz expulzate de stele de masă mică și medie, iluminate de pitica albă rămasă în centru.} \\ \midrule
\textbf{70} & \textbf{Nebuloasa Helix (\quotedblbase{}Ochiul lui Dumnezeu\textquotedblright{}, NGC 7293)} & De ce este supranumită Nebuloasa Helix (NGC 7293) \quotedblbase{}Ochiul Cosmic\textquotedblright{}? & \textbf{A: Deoarece inelul ei dublu de gaz turcoaz și roșu-auriu, cu pitica albă în mijloc, seamănă cu irisul și pupila unui ochi}. {\footnotesize Privită de-a lungul axei sale cilindrice, cavitatea albastră-turcoaz și inelul exterior roșiatic al Nebuloasei Helix creează aspectul spectaculos al unui ochi cosmic.} \\ \midrule
\textbf{71} & \textbf{Marele Nor al lui Magellan} & Ce este Marele Nor al lui Magellan și ce regiune celebră adăpostește? & \textbf{A: O galaxie-satelit a Căii Lactee (la 160.000 ani-lumină) care adăpostește uriașa Nebuloasă Tarantula}. {\footnotesize Marele Nor al lui Magellan este o galaxie-satelit a Căii Lactee ce conține Nebuloasa Tarantula, cea mai intensă maternitate stelară din Grupul Local.} \\ \midrule
\textbf{72} & \textbf{Galaxia Triunghiului (M33)} & Care sunt cele trei mari galaxii spirale care domină \quotedblbase{}Grupul Local\textquotedblright{} din care facem parte? & \textbf{A: Galaxia Andromeda (M31), Calea Lactee și Galaxia Triunghiului (M33)}. {\footnotesize Grupul Local este dominat de trei galaxii spirale: Andromeda (cea mai mare), Calea Lactee și Galaxia Triunghiului (M33).} \\ \midrule
\textbf{73} & \textbf{Galaxia Vâltoarea (M51)} & Ce particularitate celebră are Galaxia Vâltoarea (M51)? & \textbf{A: Are două brațe spirale perfecte și interacționează gravitațional cu o galaxie companion mai mică aflată la capătul unui braț}. {\footnotesize Galaxia Vâltoarea (M51) este celebră pentru brațele sale spirale grand-design și pentru dansul gravitațional cu galaxia companion NGC 5195.} \\ \midrule
\textbf{74} & \textbf{Galaxia Sombrero (M104)} & De unde provine aspectul de \quotedblbase{}pălărie Sombrero\textquotedblright{} al galaxiei M104? & \textbf{A: Este privită aproape din profil, având un bulb central uriaș de stele aurii tăiat la mijloc de un inel dens de praf cosmic}. {\footnotesize Galaxia Sombrero (M104) este văzută aproape din muchie, combinând un bulb central foarte luminos cu o bandă întunecată de praf ecuatorial.} \\ \midrule
\end{longtable}
\endgroup

\clearpage
\section{Glosar de Termeni Astronomici și Spațiali}

\begin{description}[leftmargin=1.2em, itemsep=4pt]
  \item[Albedo:] fracțiunea din lumina solară pe care un corp ceresc o reflectă înapoi în spațiu. Enceladus are un albedo de aproape $0{,}99$ (reflectă $99\%$ din lumină ca zăpada proaspătă), în timp ce asteroizii carboniferi precum Bennu reflectă doar $4\%$.
  \item[An-lumină (a.l.):] distanța parcursă de lumină în vid într-un an julian ($365{,}25\text{ zile}$), egală cu aproximativ $9.460{,}7\text{ miliarde de kilometri}$ ($63.241\text{ UA}$).
  \item[Centura Kuiper:] regiune în formă de disc situată dincolo de orbita planetei Neptun ($30\text{--}50\text{ UA}$), populată de corpuri înghețate și planete pitice precum Pluto, Haumea și Makemake.
  \item[Criovulcanism:] fenomen geologic prin care „vulcanii de gheață” de pe sateliți înghețați (precum Enceladus sau Triton) aruncă în spațiu jeturi de apă, amoniac, azot sau metan în loc de rocă topită.
  \item[Efect de seră:] procesul prin care gazele dintr-o atmosferă planetară (precum dioxidul de carbon $\text{CO}_2$ pe Venus sau metanul $\text{CH}_4$) captează căldura infraroșie, ridicând temperatura suprafeței.
  \item[Elemente kepleriene:] cei 6 parametri matematici care definesc în mod unic o orbită în spațiu: semiaxa mare ($a$), excentricitatea ($e$), înclinația ($i$), longitudinea nodului ascendent ($\Omega$), argumentul periheliului ($\omega$) și anomalia medie ($M$).
  \item[Exoplanetă:] planetă aflată în afara Sistemului Solar, care orbitează în jurul altei stele.
  \item[Gaură neagră supermasivă:] obiect cosmic cu o gravitație atât de intensă încât nici măcar lumina nu poate scăpa din interiorul \emph{orizontului evenimentelor} său; se află în nucleul majorității galaxiilor (precum \emph{Sagittarius A*} în Calea Lactee).
  \item[Magnetosferă:] scutul magnetic invizibil generat de nucleul metalic lichid al unei planete (precum Pământul sau Jupiter), care deviază particulele încărcate electric ale vântului solar.
  \item[Nebuloasă:] nor interstelar uriaș de gaz (hidrogen, heliu și gaze ionizate) și praf cosmic; poate fi o maternitate de stele (\emph{Nebuloasa Orion}) sau rămășița unei stele (\emph{Nebuloasa Inelul}, \emph{Nebuloasa Crabului}).
  \item[Norul lui Oort:] rezervor sferic uriaș de trilioane de nuclee cometare înghețate care înconjoară Sistemul Solar la distanțe de până la $\sim 1\text{ an-lumină}$.
  \item[Pitică albă:] nucleul fierbinte și extrem de dens (de mărimea Pământului, dar cu o masă comparabilă cu a Soarelui) rămas după ce o stea de masă mică sau medie își eliberează straturile exterioare.
  \item[Pitică roșie:] cel mai răspândit și mai longeviv tip de stea din Univers (clasa spectrală M, precum \emph{Proxima Centauri} și \emph{TRAPPIST-1}), mai mică și mai rece decât Soarele.
  \item[Punctul Lagrange L2:] punct de echilibru orbital situat la $1{,}5\text{ milioane de kilometri}$ în spatele Pământului (pe direcția opusă Soarelui), ideal pentru telescoapele în infraroșu precum \emph{James Webb} și \emph{Euclid}.
  \item[Rezonanță orbitală Laplace:] raportul matematic exact ($4:2:1$) dintre perioadele orbitale ale lunilor \emph{Io}, \emph{Europa} și \emph{Ganymede}, care menține interiorul lor încălzit prin forțe mareice.
  \item[Rotație sincronă (blocare mareică):] fenomenul prin care un satelit (cum este Luna față de Pământ) sau o exoplanetă apropiată se rotește în jurul axei sale exact în același timp în care parcurge o orbită, arătând mereu aceeași emisferă către corpul central.
  \item[Unitate Astronomică (UA):] distanța medie dintre Pământ și Soare ($149.597.870{,}7\text{ km}$).
  \item[Zonă locuibilă (Goldilocks):] intervalul de distanțe față de o stea în care temperatura permite existența apei lichide la suprafața unei planete.
\end{description}

\clearpage
